% Options for packages loaded elsewhere
\PassOptionsToPackage{unicode}{hyperref}
\PassOptionsToPackage{hyphens}{url}
%
\documentclass[
  11pt,
]{article}
\usepackage{amsmath,amssymb}
\usepackage{iftex}
\ifPDFTeX
  \usepackage[T1]{fontenc}
  \usepackage[utf8]{inputenc}
  \usepackage{textcomp} % provide euro and other symbols
\else % if luatex or xetex
  \usepackage{unicode-math} % this also loads fontspec
  \defaultfontfeatures{Scale=MatchLowercase}
  \defaultfontfeatures[\rmfamily]{Ligatures=TeX,Scale=1}
\fi
\usepackage{lmodern}
\ifPDFTeX\else
  % xetex/luatex font selection
\fi
% Use upquote if available, for straight quotes in verbatim environments
\IfFileExists{upquote.sty}{\usepackage{upquote}}{}
\IfFileExists{microtype.sty}{% use microtype if available
  \usepackage[]{microtype}
  \UseMicrotypeSet[protrusion]{basicmath} % disable protrusion for tt fonts
}{}
\makeatletter
\@ifundefined{KOMAClassName}{% if non-KOMA class
  \IfFileExists{parskip.sty}{%
    \usepackage{parskip}
  }{% else
    \setlength{\parindent}{0pt}
    \setlength{\parskip}{6pt plus 2pt minus 1pt}}
}{% if KOMA class
  \KOMAoptions{parskip=half}}
\makeatother
\usepackage{xcolor}
\usepackage[margin=1in]{geometry}
\setlength{\emergencystretch}{3em} % prevent overfull lines
\providecommand{\tightlist}{%
  \setlength{\itemsep}{0pt}\setlength{\parskip}{0pt}}
\setcounter{secnumdepth}{-\maxdimen} % remove section numbering
\ifLuaTeX
  \usepackage{selnolig}  % disable illegal ligatures
\fi
\usepackage{bookmark}
\IfFileExists{xurl.sty}{\usepackage{xurl}}{} % add URL line breaks if available
\urlstyle{same}
\hypersetup{
  pdftitle={Bounded local moment measures and the geometric curve reference case},
  pdfauthor={Fresh partial reconstruction of original P12.4},
  hidelinks,
  pdfcreator={LaTeX via pandoc}}

\title{Bounded local moment measures and the geometric curve reference
case}
\author{Fresh partial reconstruction of original P12.4}
\date{8 October 2026}

\begin{document}
\maketitle

{
\setcounter{tocdepth}{3}
\tableofcontents
}
\section{Bounded local moment measures and the geometric curve reference
case}\label{bounded-local-moment-measures-and-the-geometric-curve-reference-case}

Partial progress on original P12.4. This manuscript does \textbf{not}
claim a complete singular/nonnormal geometric Monge--Ampere
construction. It proves an algebraic bounded moment extension and an
actual geometric zero-potential curve reference case. Independent
higher-dimensional Cartier intersection comparison and open-set locality
remain unproved. The final candidate flag is false.

The mathematical arguments, published-source applications, finite
observations, and artifact integrity checks are distinct. The only
incoming grant is qualified G331. Its complete source, mandatory
annotation, ledger, and exact scope are retained unchanged locally. No
supplied executable or unreviewed route is a theorem premise.

\section{n1. Domain, reference data, and the retained
grant}\label{n1.-domain-reference-data-and-the-retained-grant}

This manuscript is a partial result for P12.4, not a solution claim. The
only supplied result used is qualified G331, including its entire
source, mandatory annotation, source ledger, and exact scope. ROOT338 is
an unreviewed approach, not a premise. No envelope, intersection,
product, energy, locality, or arbitrary-model cofinality theorem is
imported from G331.

Let \((R,\mathfrak m)\) be excellent, reduced, local, and of positive
dimension \(n\). Put \(X_{\mathfrak a}=\{v:v(\mathfrak a)=1\}\) for a
fixed proper \(\mathfrak m\)-primary ideal \(\mathfrak a\). A point is a
centred real semivaluation on \(R\) with \(0<v(\mathfrak m)<\infty\);
its prime support may be nonminimal. Every such support is retained,
including supports on the conductor quotient. Ideal values mean minima
on a finite generating set. In particular primary ideal values are
finite, positive, continuous, and bounded on \(X_{\mathfrak a}\). Units
have value zero. G331 gives the compact Hausdorff boundary and the exact
primary change \(T(v)=v/v(\mathfrak b)\).

The geometric core was defined before ideal closure: actual negative
effective finite REAL Cartier presentations on actual primary blowups,
with nonnegative degree on every contracted proper integral curve. G331
proves that its uniform closure \(C=C_{\rm geom}\) equals the uniform
closure of finite nonnegative real combinations of \(-v(I)\), where
\(I\) is primary or the unit ideal. Its completion \(D=D_{\rm geom}\)
consists of infima of ACTUAL nonempty decreasing \(C\)-nets, with the
identically \(-\infty\) function allowed. We use only this completion,
not closure under arbitrary infima, increasing limits, or flattened
nested nets.

The product-input domain in this manuscript is
\[D^{\rm bd}_{\mathfrak a}=\{u\in D:-M\le u\le0\text{ for some finite }M\}.\]
This includes every bounded real, discontinuous completed input, zero,
and permitted mixed-zero branch patterns. It excludes every input having
a \(-\infty\) value, including bottom. Every continuous input in \(C\)
is bounded. Existence of a bounded decreasing representation needs no
truncation: if \(g_\alpha\downarrow u\) and \(u\ge-M\), then
\(g_\alpha\ge u\ge-M\) and \(g_\alpha\le0\).

Use \(p=n-1\) potential entries; a test supplies the nth intersection
entry. For non-equidimensional \(R\), all algebraic moments below use
the degree-\(n\) part of the primary length polynomial.
Lower-dimensional components remain in the boundary and domain, but
contribute zero in this chosen convention. A separate lower-dimensional
product would have a different arity; none is silently added. In
dimension one there are zero potentials and a reference measure,
constructed geometrically in n5. For dimension zero the reduced local
ring is a field, the centred boundary is empty, and the measure is the
unique zero measure; there is no negative-arity assertion.

G331 also supplies two specific functional facts used below. Rational
profiles \(-v(I)/k\) with \(k\in\mathbb N_{>0}\) (unit allowed) are
uniformly dense in \(C\). The rational difference span of ideal gauges
is a lattice, contains rational constants since \(v(\mathfrak a)=1\),
and is uniformly dense in \(C(X_{\mathfrak a},\mathbb R)\). To recall
the mechanism, a common denominator turns any span element into
\((v(J)-v(K))/k\), and \[\min(v(I)-v(J),v(K)-v(H))=v(IH+KJ)-v(JH).\] The
bounded truncations \(v((f)+\mathfrak a^N)=\min(v(f),N)\) separate all
semivaluations, including infinite coordinates. Uniform closure of this
rational lattice is a real vector lattice; the compact lattice
approximation argument yields all real continuous tests. This is density
on the entire compact boundary, not a finite grid.

For a finite normalization \(R\hookrightarrow S\), G331 retains the
disjoint local normal centres above \(\mathfrak m\) and their FULL
restriction fibres. At finite conductor value the fibre is a singleton.
At infinite conductor value the identification is exactly equality on
the conductor quotient. Continuous and completed geometric functions
descend precisely when constant on these whole fibres. We do not select
one lift or discard an infinity stratum. Normalization in n4 is applied
only to actual pullbacks of downstairs functions.

For nonreduced germs the present data are those of the reduction and the
reduced gauge. No nilpotent generic length is reconstructed; n6 proves
an explicit obstruction to doing so from these data alone. P12.5 energy
and general finite-map geometric operators in P12.6 are outside the
results here.

The paper and local notes give mathematical draft arguments; hashing,
compilation, and finite diagnostics are separate evidence. The entire
G331 source/annotation/ledger are retained under inputs/mathematical.
Rerun data retention: \texttt{python\ audit\_inputs.py}. New source
applications are recorded in source-application-ledger.md.

\section{n2. An unconditional primary mixed-multiplicity
kernel}\label{n2.-an-unconditional-primary-mixed-multiplicity-kernel}

Use the domain and full-boundary conventions of n1. This note proves the
algebraic kernel needed for moment measures, with arbitrary
characteristic, residue field, non-Cohen--Macaulay rings, and
non-equidimensional reduced components allowed. It proves no geometric
Monge--Ampere operator.

Here are the precise published applications. In the CURRENT July 2024
author version of Swanson--Huneke, Theorem 17.4.2 says that for a
Noetherian local ring, finite module \(M\), and primary ideals
\(I_1,\ldots,I_s\), the function
\(\ell_R(M/I_1^{t_1}\cdots I_s^{t_s}M)\) is a rational polynomial of
total degree \(\dim M\) for all sufficiently large independent integer
\(t_i\). Definition 17.4.3 takes its degree-\(\dim R\) coefficients,
using factorial denominators, as mixed multiplicities; they are zero if
\(\dim M<\dim R\). These statements and their full immediate proof were
read. We use this degree convention even when components have different
dimensions.

The complete 17.2.2 proof gives a nonempty Zariski open set of choices
for a superficial element in \(I_1/\mathfrak m I_1\), with finitely many
prescribed primes avoidable when \(I_1\) is not contained in them.
Complete 17.2.4 and 17.4.6 show that a superficial element \(x\in I_1\)
avoiding the minimal primes reduces a mixed multiplicity having at least
one \(I_1\) entry to the corresponding mixed multiplicity over \(R/xR\)
with that entry removed, when \(\dim R>1\). No regular sequence or
Cohen--Macaulay hypothesis is required. The finite annihilator
correction in dimension one is part of the source, but we do not need
that formula to compare the last ideal.

Finite residue fields cause no restriction. Pass to
\(R(Z)=R[Z]_{\mathfrak m R[Z]}\) solely for these algebraic proofs.
Lemma 8.4.2 gives faithful flatness, equal dimension, preservation of
primary ideals and lengths. Explicitly, a composition series of any
finite-length \(R\)-module tensors to a series with factors
\(\kappa(\mathfrak m)(Z)\), so its length is unchanged. Thus every
multigraded length polynomial and its coefficients are unchanged. No
identification of boundaries across this transcendental extension is
used.

Write \(e_R(I_1,\ldots,I_n)\) for the coefficient with one entry of each
ideal. Product and power linearity follow directly from the polynomial:
replace the variables for \(I\) and \(J\) by the same variable to
represent \((IJ)^t\). The coefficient with exponent one in that variable
and one in the other \(n-1\) variables receives exactly the terms with
exponent one in \(I\) or in \(J\). Hence
\[e_R(IJ,I_2,\ldots,I_n)=e_R(I,I_2,\ldots,I_n)+e_R(J,I_2,\ldots,I_n),\]
and \(e_R(I^r,I_2,\ldots,I_n)=r e_R(I,I_2,\ldots,I_n)\). Relabelling
gives symmetry. If an ideal is the unit ideal, its length polynomial is
independent of that variable, so the corresponding mixed multiplicity is
zero. This convention is compatible with the identities.

We give the single-slot inclusion proof, rather than trying to compare
individual coefficients of unequal multivariate polynomials. Suppose
\(J\subset K\) are primary and the other \(n-1\) ideals are fixed. Over
an infinite residue field, choose in the second slot an element
superficial for BOTH tuples \((I_2,I_3,\ldots,I_n,J)\) and
\((I_2,I_3,\ldots,I_n,K)\), and avoiding all minimal primes. The two
allowed sets are nonempty opens in the same finite-dimensional vector
space \(I_2/\mathfrak m I_2\), obtained by removing finitely many proper
subspaces in the 17.2.2 proof. Their intersection still has a rational
point over an infinite field. The avoidance conditions can be included
in that finite union. Each quotient has dimension one less: avoiding the
minimal primes gives the upper bound, and lifting a system of parameters
from the quotient and using the height theorem gives the lower bound.
The quotient images of \(J,K\) remain primary and included. Apply 17.4.6
to both tuples. Repeat in the remaining fixed slots, choosing
simultaneous superficial elements afresh in each quotient. In dimension
one, \(J^N\subset K^N\) gives
\[\ell(R/J^N)\ge\ell(R/K^N),\qquad e(J)\ge e(K)\] by dividing by \(N\)
and using the one-variable Hilbert--Samuel polynomial. Thus, in every
dimension,
\[J\subset K\quad\Longrightarrow\quad e_R(J,I_2,\ldots,I_n)\ge e_R(K,I_2,\ldots,I_n).\]
The same reduction to a one-dimensional primary multiplicity proves
positivity for primary slots. The positivity of the last one-variable
leading coefficient is Theorem 11.1.3 and Definition 11.1.5, whose
immediate proof/context were read. Unit slots give zero. These arguments
also work for the nonreduced quotients created by superficial elements.

Integral-closure invariance has a separate proof. If
\(K\subset\overline J\), put \(L=J+K\). Each degree-one generator \(xt\)
of \(R[Lt]\) is integral over \(R[Jt]\) by its homogeneous
ideal-integrality equation. Thus \(R[Lt]\) is a finite graded
\(R[Jt]\)-module. A finite homogeneous generating list, with degrees at
most \(h\), gives \(L^{t+h}\subset J^t\) for all \(t\ge1\): its degree
\(t+h\) coefficients lie in \(J^{t+h-d}\subset J^t\). Also
\(J^t\subset L^t\). After multiplying by any fixed products
\(A=I_2^{t_2}\cdots I_s^{t_s}\), this gives
\[\ell(R/L^tA)\le\ell(R/J^tA)\le\ell(R/L^{t+h}A).\] On every ray of
positive integral exponents, divide by the common large parameter to the
power \(n\). The two outside polynomials have the same degree-\(n\)
homogeneous part; the fixed shift \(h\) contributes only smaller
degrees. Thus the degree-\(n\) parts for \(J\) and \(L\) agree on every
positive rational ray and therefore as polynomials. All mixed
coefficients agree. If \(J,K\) have equal integral closure, apply this
twice to \(J+K\). This proves invariance without transferring
coefficient inequalities from a length inequality. The finite graded
argument is proved here and does not require finite normalization of a
Rees algebra.

Finally, the order used in these moments is the order of functions on
ALL \(X_{\mathfrak a}\). For primary or unit \(J,K\),
\[v(K)\ge v(J)\text{ on }X_{\mathfrak a}\quad\Longleftrightarrow\quad K\subset\overline J.\]
For proper primary \(J\), apply the complete Swanson--Huneke 6.8.4
centred rank-one DVR criterion to each domain \(R/P\) for
\(P\in\operatorname{Min}R\). Each quotient has positive dimension: in a
positive-dimensional reduced local ring the maximal ideal cannot be
minimal. Normalize these DVRs by the nonzero primary image of
\(\mathfrak a\); they are actual points of \(X_{\mathfrak a}\) with
support \(P\). Then 6.8.3 (or 1.1.5 componentwise integral dependence)
assembles the conclusion. Conversely a monic ideal-integrality equation
and the unique-smallest-term argument in the quotient by any
semivaluation support give the inequality at that support; primary
values are finite. For \(J=R\) the conclusion is automatic. For \(K=R\)
and proper \(J\) the inequality is impossible on the nonempty section.
These cover unit cases. Nonminimal supports are not removed merely
because the criterion can be tested on a subset.

For rational profiles \(u_i=-v(I_i)/k_i\) and a primary ideal test
\(v(J)\), define
\[B_{\boldsymbol u}(v(J))=\frac{e_R(I_1,\ldots,I_p,J)}{k_1\cdots k_p}.\]
For \(p=0\) this means \(e_R(J)\). A zero input may be represented by
the unit ideal and gives zero if \(p\ge1\). Equal rational profiles
have, after integer powers, equal primary ideal values, hence equal
integral closures; invariance and power linearity make the definition
independent of presentation. Any rational linear relation among ideal
tests becomes \(v(A)=v(B)\) after synchronizing denominators and
multiplying ideals, so the same argument makes the test functional
well-defined. If a rational test function \(v(A)-v(B)\) is nonnegative,
the criterion gives \(A\subset\overline B\). Since \(A+B\) has the same
integral closure as \(B\), inclusion monotonicity and invariance give
\(e(\boldsymbol I,A)\ge e(\boldsymbol I,B)\). Thus the functional is
positive on the whole rational test lattice.

Input order also follows: \(u\le w\) implies the moments for \(u\) are
at least those for \(w\) in that slot. For \(u=-v(I)/k,w=-v(J)/l\), use
\(I^l\subset\overline{J^k}\) and divide by \(kl\). All slots are
symmetric and additive on addition of negative rational profiles. These
are the algebraic kernel and order laws, not Cartier intersection
comparison or locality.

Source bodies/hashes/URLs and the exact units used are retained in
inputs/primary, reports/input-hashes.json, and
source-application-ledger.md. Rerun: \texttt{python\ audit\_inputs.py};
\texttt{python\ build\_delivery.py}. No program establishes these
universal proofs.

\section{n3. Bounded real and actual decreasing-net moment
measures}\label{n3.-bounded-real-and-actual-decreasing-net-moment-measures}

Use n1 and the proved algebraic kernel n2. This note constructs finite
positive Radon MOMENT measures \(\nu_{\mathfrak a}(u_1,\ldots,u_p)\) on
the entire boundary for every tuple in \((D^{\rm bd}_{\mathfrak a})^p\).
The name deliberately does not assert the geometric
intersection/locality properties still missing in n7. The proof covers
arbitrary real continuous and bounded completed functions, not just
rational profiles.

For a rational tuple from n2 its positive functional on the rational
test lattice has mass
\[B_{\boldsymbol u}(1)=B_{\boldsymbol u}(v(\mathfrak a))=\frac{e_R(I_1,\ldots,I_p,\mathfrak a)}{\prod k_i}.\]
If \(q>\|f\|_\infty\) is rational, \(q\pm f\) are nonnegative lattice
tests. Positivity gives \(|B(f)|\le qB(1)\), hence
\(|B(f)|\le\|f\|_\infty B(1)\). Uniform density from n1 extends this
rational functional uniquely to a bounded real-linear functional on all
real continuous functions: real linearity follows by rational
approximation of scalars and norm continuity. Positivity persists by
approximating a nonnegative continuous function and adding a vanishing
positive rational constant to the approximation. The compact-Hausdorff
Riesz correspondence identifies it with a finite positive regular Borel
(Radon) measure. Equivalently, in the convention explicitly used in BFJ
Section 4.2, a Radon measure is its positive continuous test functional.
The Riesz correspondence is ordinary functional background at exactly
compact Hausdorff hypotheses, not an extra geometric positivity theorem.

There is a uniform bound on gauge moments. If \(-M\le u_i\le0\) for
every rational entry, order comparison with the constant \(-M\) (first
take rational constants just above \(M\)) gives
\[0\le B_{\boldsymbol u}(v(J))\le M^p e_R(\mathfrak a^{[p]},J),\qquad \nu_{\mathfrak a}(\boldsymbol u)(X_{\mathfrak a})\le M^p e_R(\mathfrak a).\]
Here brackets mean repeated ENTRIES, not ideal exponentiation;
\(e_R(\mathfrak a)=e_R(\mathfrak a^{[n]})\). For \(p=0\) the factor is
one. No order of measures, as opposed to order on ideal-gauge moments,
is asserted.

The decisive error estimate uses legitimate negative constants in \(C\).
For rational tuples \(\boldsymbol u,\boldsymbol w\) in \([-M,0]\) with
coordinatewise sup-norm difference at most \(\varepsilon\), use
\(u_i-\varepsilon\le w_i\) and the reverse inequalities, with a rational
slightly larger error if necessary. Expanding the additive kernel after
adding \(-\varepsilon=-\varepsilon v(\mathfrak a)\) in each slot gives
\[\left|B_{\boldsymbol u}(v(J))-B_{\boldsymbol w}(v(J))\right|\le [(M+\varepsilon)^p-M^p]e_R(\mathfrak a^{[p]},J).\tag{1}\]
Indeed each nonempty subset of shifted slots contributes a nonnegative
term bounded by \(\varepsilon^{|S|}M^{p-|S|}e_R(\mathfrak a^{[p]},J)\);
summing them is the binomial expression. For \(p=0\) the expression is
zero.

Now take arbitrary REAL \(C\) inputs and approximate each uniformly by
rational profiles using n1. The approximants may lie in \([-M-1,0]\),
which supplies a common bound. Estimate (1) makes every ideal-test
moment Cauchy and independent of approximation. Every rational lattice
test is a difference of such gauges, so its functional values converge.
Positivity and the mass bound pass to the limit, hence give a measure by
the preceding extension argument. Finite uniform density then proves
weak convergence on every continuous test. The same estimate with
arbitrarily slightly larger bounds/errors passes to the real \(C\)
inputs. Symmetry, nonnegative real homogeneity, additivity within the
negative cone, and order on ideal-gauge moments pass to the limit. For
additivity use rational approximations of each summand and their sum,
and for real homogeneity approximate the scalar rationally. These are
properties of the proved moment functional; the proof does not cite BFJ
4.4 outside its smooth range.

For each \(u_i\in D^{\rm bd}_{\mathfrak a}\) choose an ACTUAL decreasing
\(C\)-net \(g_{i,\alpha_i}\downarrow u_i\). By n1 all its terms have the
bound \([-M,0]\) if the tuple has that bound. The finite Cartesian
product of the directed index sets represents the mixed tuple. For every
ideal test \(v(J)\), its moments increase under refinement and stay
bounded by the preceding estimate. Define \(B_{\boldsymbol u}(v(J))\) as
their finite supremum. This is an increasing limit of bounded SCALARS,
not an increasing limit in the function class. Values for difference
tests converge by subtraction. Positive test functionals and their
common mass bound pass to the limit on the rational lattice and extend
to a finite positive measure as before.

We prove representation independence without a flattening assumption.
Let \(g_{i,\alpha_i}\) and \(f_{i,\beta_i}\) represent the same bounded
tuple. Fix a tuple of indices \(\boldsymbol\beta\) and
\(\varepsilon>0\). Since \(u_i\le f_{i,\beta_i}\), the directed open
sets \[\{g_{i,\alpha_i}<f_{i,\beta_i}+\varepsilon\}\] cover the compact
space. A finite subcover and directedness select one \(\alpha_i\) for
which the inequality holds everywhere. Do this for each of the finitely
many slots. Then \(g_{i,\alpha_i}-\varepsilon\le f_{i,\beta_i}\). Moment
order and the real \(C\) negative-shift expansion give
\[B_{\boldsymbol f_{\beta}}(v(J))\le B_{\boldsymbol g_{\alpha}-\varepsilon}(v(J))\le \sup_{\boldsymbol\alpha}B_{\boldsymbol g_{\alpha}}(v(J))+[(M+\varepsilon)^p-M^p]e_R(\mathfrak a^{[p]},J).\]
Take the supremum over \(\boldsymbol\beta\), send \(\varepsilon\) to
zero, and reverse the roles. Equality of every gauge moment, density,
and the mass bound give equality of the measures. The same argument, now
with limit tuples satisfying \(u_i\le w_i\) and fixed continuous terms
majorizing \(w_i\), proves order on ideal-gauge moments throughout the
bounded \(D\) domain.

Addition of two actual decreasing \(C\)-nets, indexed by their finite
product, gives an actual decreasing \(C\)-net to their sum. Thus the sum
belongs to \(D\) by its definition, with a finite bound. Taking scalar
moment limits in the \(C\) additivity identity proves additivity on
bounded \(D\) entries. Scaling is handled the same way; scalar zero
produces zero. Symmetry is immediate. In particular the
negative-constant shift expansion and (1) are now valid on bounded \(D\)
inputs. No closure under arbitrary \(D\) infima or arbitrary nested-net
flattening was invoked.

The useful exact formula, for every ideal \(J\), is
\[B_{\boldsymbol u}(v(J))=\sup\{B_{\boldsymbol g}(v(J)):g_i\in C,\ g_i\ge u_i\text{ for all }i\}.\tag{2}\]
Every such \(g_i\) lies in the common bound \([-M,0]\). The order just
proved gives the upper bound; the terms of any representing nets give
the reverse bound. This is a formula for scalar moments, not a new
definition of \(D\) by all majorants.

Here is the precise descending-net theorem. Let \(\lambda\) range over
one directed set, let \(u_{i,\lambda}\in D^{\rm bd}_{\mathfrak a}\)
decrease pointwise to a SPECIFIED \(u_i\in D^{\rm bd}_{\mathfrak a}\),
and suppose all entries and limits have a common finite bound
\([-M,0]\). Then
\[\nu_{\mathfrak a}(u_{1,\lambda},\ldots,u_{p,\lambda})\ \longrightarrow\ \nu_{\mathfrak a}(u_1,\ldots,u_p)\]
weakly, for the net. The membership of the limit is an actual
hypothesis; no new infimum closure is being claimed. To prove it, order
gives a bound above by the target gauge moment. Fix a continuous
majorant tuple \(g_i\ge u_i\). Every \(D\) function is upper
semicontinuous, so \(\{u_{i,\lambda}<g_i+\varepsilon\}\) is open. The
decreasing net gives a directed cover; compactness selects one common
\(\lambda\) for all slots with \(u_{i,\lambda}-\varepsilon\le g_i\).
Thus
\[B_{\boldsymbol u_\lambda}(v(J))\ge B_{\boldsymbol g}(v(J))-[(M+\varepsilon)^p-M^p]e_R(\mathfrak a^{[p]},J).\]
First send \(\varepsilon\) to zero, then take the supremum in (2). The
gauge moments converge to the target moments. They determine convergence
for the rational difference lattice; density and the uniform mass bound
extend it to every continuous test. This proves the actual weak
descending-net theorem for this moment construction. It is not
convergence in total variation, and it is not locality.

All-zero inputs yield zero for \(p\ge1\), and any single globally zero
entry annihilates the mixed measure. For mixed-zero patterns on
normalization branches, the whole tuple is still in the stated domain if
it belongs to G331's class; this construction makes no strict-negativity
exclusion. It is not extended here to unbounded or bottom inputs.

Artifacts: n3.md, argument.md, P12.4-paper.tex, and native-notes.json.
Rerun: \texttt{python\ build\_delivery.py}. The conclusion is a unique
bounded moment extension with the displayed laws. Independent Cartier
comparison and actual open-set locality remain separate obligations.

\section{n4. Exact gauge, normalization, and published smooth comparison
for
moments}\label{n4.-exact-gauge-normalization-and-published-smooth-comparison-for-moments}

Use G331 and n1--n3. This note concerns the moment measures proved
there. The rules are genuine identities for those measures, but do not
fill the missing singular Cartier comparison or locality theorem.

For two proper PRIMARY gauges \(\mathfrak a,\mathfrak b\), let
\(d(v)=v(\mathfrak b)\) on \(X_{\mathfrak a}\) and \(T(v)=v/d(v)\). The
function \(d\) is continuous and bounded above and away from zero. For
each degree-one homogeneous input set
\(u_{i,\mathfrak b}(T(v))=u_{i,\mathfrak a}(v)/d(v)\). G331 proves that
this transports both \(C\) and actual \(D\) representations. Bounds
remain finite, and uniform convergence transports because \(1/d\) is
bounded. Then the EXACT measure change for the moment construction is
\[\nu_{\mathfrak b}(\boldsymbol u_{\mathfrak b})=T_*\big[d\,\nu_{\mathfrak a}(\boldsymbol u_{\mathfrak a})\big].\tag{3}\]
For a rational ideal tuple and a gauge test \(v_{\mathfrak b}(J)\), the
right side pairs to
\[\int d(v)\frac{v(J)}{d(v)}\,d\nu_{\mathfrak a}=\frac{e_R(I_1,\ldots,I_p,J)}{\prod k_i},\]
which is the left moment. Density of the test lattice gives equality of
measures. Uniform and actual descending continuity from n3 then prove
(3) for all bounded inputs in \(D\). Multiplication of tests by the
fixed continuous bounded \(d\) commutes with those weak limits. This is
one factor of \(d\), not \(d^n\) or \(d^{-(n-1)}\). For
\(\mathfrak b=\mathfrak a^r\), the same ray has \(d=r\), the potentials
on its new section divide by \(r\), and the transported measure
multiplies by \(r\). The cocycle follows by substitution. There is no
nonprimary denominator anywhere.

Let \(S\) be the finite normalization of reduced excellent \(R\), and
let \(q\) run over all its local normal centres above \(\mathfrak m\).
Put \(S_q\) for the corresponding normal local domain,
\(f_q=[\kappa(q):\kappa(\mathfrak m)]\), and
\(\rho_q:X_{\mathfrak a S_q}\to X_{\mathfrak a}\) for restriction. Use
only centres with \(\dim S_q=n\) in the degree-\(n\) formula:
\[\nu_{\mathfrak a,R}(\boldsymbol u)=\sum_{q:\dim S_q=n} f_q\,\rho_{q,*}\nu_{\mathfrak a S_q,S_q}(\rho_q^*\boldsymbol u).\tag{4}\]
All other centres and all semivaluation supports are retained; a centre
of smaller dimension has zero top-degree contribution in this
convention, rather than a fictitious degree assigned to it.

Here is the algebraic proof. Complete Swanson--Huneke 17.4.4 gives
additivity of mixed multiplicity on exact sequences of finite modules,
and complete 17.4.8 gives the sum over maximal-dimensional minimal
primes with coefficient \(\ell(M_P)\). Applied to \(M=R,S\), the
coefficients at every minimal \(P\) of \(R\) are both one: a reduced
ring localized at a minimal prime is its component field, and
normalization is generically the same product of fields. Hence
\(e_R(\boldsymbol I;R)=e_R(\boldsymbol I;S)\). Alternatively \(S/R\) is
annihilated by a regular conductor element and has dimension less than
\(n\), so additivity gives the same identity. For each product
\(A=I_1^{t_1}\cdots I_s^{t_s}\), the finite Artinian quotient \(S/AS\)
splits into its local factors. A finite-length \(S_q\) composition
series has \(R\)-length \(f_q\) per factor, so
\[\ell_R(S/AS)=\sum_q f_q\ell_{S_q}(S_q/AS_q).\] Compare degree-\(n\)
homogeneous coefficients using 17.4.2. Terms with \(\dim S_q<n\) are
zero; this gives
\[e_R(I_1,\ldots,I_n)=\sum_{q:\dim S_q=n} f_q e_{S_q}(I_1S_q,\ldots,I_nS_q).\]
Formula (4) follows on rational tuples from the gauge moments and test
density, and on \(C\) and bounded \(D\) tuples by the proved
continuities. Pullback of each chosen decreasing representation is an
actual bounded decreasing representation on each centre, by G331. There
are finitely many centres, so summation commutes with the limits.

This proof checks actual residue degrees, including inseparable ones. It
uses generic identity of normalization on reduced component fields, not
an invented vertical degree and not a general finite-map trace theorem.
On the conductor-infinite stratum every normalization point is still
present. Formula (4) pushes all of them down, adding any masses whose
fibres are identified. The downstairs inputs restrict with identical
values on the entire fibre, by their actual pullback, and the formula is
unaffected by descriptions of that fibre. No assertion of automatic
descent of arbitrary upstairs operator data is made. General finite
maps/P12.6 remain separate.

At the EXACT published BFJ range there is an additional comparison. Let
the germ be at a smooth CLOSED point of a smooth affine variety over an
algebraically closed field of characteristic zero, of dimension
\(n\ge2\), with \(\mathfrak a=\mathfrak m\). G331 identifies \(C,D\)
with the full BFJ classes on all semivaluations, with bottom adjoined;
our bounded domain excludes bottom. BFJ's setting was read at its actual
hypotheses. Its complete Section 4 was read, including proofs of
4.1--4.9. It defines the independent Cartier intersection, gives
\[-\langle Z(I_1),\ldots,Z(I_n)\rangle=e_R(I_1,\ldots,I_n),\] and in
Proposition 4.6 characterizes its mixed measure by pairing psh test
functions with that intersection. Proposition 4.4 proves symmetry,
additivity, order, and continuity of nef intersections; its proof and
Remark 4.5 explicitly use smooth multiplier approximation for
additivity. Proposition 4.7 gives the ideal moments, 4.8
decreasing-sequence weak continuity, and 4.9 the model-atom formula.
These are NOT applied at a singular/positive-characteristic germ.

The polynomial mixed multiplicity in n2 agrees here with BFJ's
generic-length convention: choose a superficial parameter sequence in
the given ideals, as in complete SH 17.4.7. The local ring is regular,
hence Cohen--Macaulay, so the successive parameter elements are
nonzerodivisors. Repeated SH 17.4.6 reduces to its dimension-one
formula, whose annihilator term is zero. The coefficient is therefore
the length of the generic parameter quotient, exactly the convention in
BFJ (4.3). This equality of conventions is used only in this regular
range.

On rational ideal tuples, BFJ 4.7 consequently gives exactly n2's gauge
moments, so test density identifies its measure with
\(\nu_{\mathfrak m}\). For \(C\) tuples use uniform rational
approximation. The order of BFJ intersections and their additivity give
the same negative-shift estimate (1), hence uniform weak convergence,
and equality follows. Finally let \(g_{i,\alpha_i}\downarrow u_i\) be
any actual bounded decreasing \(C\) representations of bounded BFJ
inputs. Their associated Weil coefficients decrease to those of \(u_i\)
at every divisorial valuation. Proposition 4.4's continuity along
decreasing families applies to this finite product net, with the fixed
Cartier test \(Z(J)\). Its proof gives families, not just sequences: the
intersection is upper semicontinuous in the coefficient topology and
monotone, and a coordinatewise decreasing net converges in that
topology. Thus the gauge moments converge to those of the BFJ measure.
N3 and density identify \(\nu\) with BFJ on EVERY bounded completed
input and ALL semivaluations. For another primary gauge, equation (3)
transports this comparison; it does not attribute an unstated
arbitrary-gauge theorem to BFJ.

BFJ's exact smooth intersection comparison is inherited in this range
because the moment measure was identified with its published operator.
Its Section 4 contains no general open-set locality statement used here.
We make no new claim that this argument proves locality, even at this
published range. Curve/DVR and empty dimension-zero conventions are
n5/n1.

Artifacts and rerun: source-application-ledger.md;
\texttt{python\ audit\_inputs.py}; \texttt{python\ build\_delivery.py}.
The full BFJ PDF/text and current Swanson--Huneke bodies/hashes/URLs are
retained locally. Their ordinary recursive historical dependencies are
disclosed; execution is not a proof.

\section{n5. Actual geometric curve reference measure and dimension-zero
convention}\label{n5.-actual-geometric-curve-reference-measure-and-dimension-zero-convention}

Use qualified G331 and n1--n4. Here there IS an actual independent
Cartier intersection comparison and locality theorem, in the exact range
of every excellent reduced local CURVE, with arbitrary characteristic
and residue field. Arity is zero, so this is a reference measure and not
a positive-arity higher-dimensional product.

Let \(S\) be the finite normalization, with local DVR centres \(S_q\),
residue degrees \(f_q=[\kappa(q):\kappa(\mathfrak m)]\), and orders
\(a_q=\operatorname{ord}_q(\mathfrak a)>0\). Every centred semivaluation
of a DVR is a positive multiple of its order: write a nonzero element as
a power of a uniformizer times a unit, whose value is zero. A support
cannot be its maximal ideal, because that would make the boundary ideal
infinite. There is no other nonzero prime support in a DVR. Thus G331's
normalization presentation identifies the entire downstairs boundary
with the finite discrete set \[v_q=\operatorname{ord}_q/a_q.\] A regular
conductor element has finite value on every branch, so ALL fibres in
this curve case are singletons. These are local centres of
normalization, not merely minimal components.

Define the reference measure directly from Cartier order on normalized
curves: \[\mu_{\mathfrak a,R}=\sum_q f_q a_q\,\delta_{v_q}.\tag{5}\]
This definition uses actual residue field degrees; there is no generic
degree assigned to a vertical component. Every coefficient is positive
and finite. The construction is geometric before any ideal-moment
comparison.

We prove comparison with an arbitrary actual primary blowup model
\(Y=\operatorname{Bl}_I R\) and every exceptional REAL Cartier divisor
\(F\) on it. Since \(I S\) is invertible on the semilocal normal curve,
the universal property gives a morphism \(h:\operatorname{Spec}S\to Y\).
It is finite: its graph is closed in \(\operatorname{Spec}S\times_R Y\)
because \(Y\) is separated over \(R\), and the projection of this
product to \(Y\) is finite as a base change of the finite normalization.
Its image contains the generic point of every component of \(Y\), which
is reduced and generically birational on all components by the primary
blowup construction in G331. The finite image is closed, hence the map
is surjective. It is the normalization of \(Y\): locally it is a finite
normal birational extension with the same total quotient fields. This
special curve argument establishes domination of these actual primary
models. It makes no arbitrary higher-dimensional proper-model cofinality
assertion.

For an effective integral exceptional Cartier divisor \(A\), its
independent one-dimensional intersection degree is the proper zero-cycle
degree
\[\deg_R[A]=\sum_{y\in\operatorname{Supp}A}\ell_{\mathcal O_{Y,y}}(\mathcal O_{Y,y}/(g_y))\,[\kappa(y):\kappa(\mathfrak m)],\]
where \(g_y\) is a regular local equation. The support is a finite
subset of the closed fibre: the finite normalization surjects onto it,
and its fibre above \(\mathfrak m\) is finite. Each relevant local ring
is reduced of dimension one, and \(A\) cuts a finite-length quotient.
Its residue field is finite over the base residue field because it is a
closed point of the proper zero-dimensional fibre. This is the usual
local Cartier zero-cycle, not a degree defined from valuative moments.

Here is a direct normalization proof of that degree. Put
\(T=\mathcal O_{Y,y}\) and let \(N\) be its finite semilocal
normalization. The cokernel \(C=N/T\) has finite length, since
normalization is generically identical on a reduced curve.
Multiplication by the regular element \(g=g_y\) in
\(0\to T\to N\to C\to0\) gives
\[0\to C[g]\to T/gT\to N/gN\to C/gC\to0.\] For the finite-length module
\(C\), its multiplication exact sequence gives
\(\ell(C[g])=\ell(C/gC)\). Therefore \(\ell_T(T/gT)=\ell_T(N/gN)\). The
DVR factors of \(N\) give
\[\ell_T(N/gN)=\sum_{q\mapsto y}[\kappa(q):\kappa(y)]\operatorname{ord}_q(g).\]
Multiplying by \([\kappa(y):\kappa(\mathfrak m)]\) and summing proves
\[\deg_R[A]=\sum_q f_q\operatorname{ord}_q(h^*A).\tag{6}\] This proof
works without normality or Cohen--Macaulay assumptions on \(Y\). It uses
local actual Cartier equations and the finite normalization, not a
replacement by nonnormal Weil coefficients.

Every integral exceptional Cartier \(F\) on a primary model can be
written as a difference of effective exceptional Cartier divisors by the
actual Cartier order-unit assertion of G331; extend (6) by differences
and then by finite real linear combinations. Relations cause no
ambiguity, since each local order is a homomorphism on actual Cartier
divisors. The actual intersection degree of \(F\) is therefore
\[\deg_R(F)=\sum_q f_q\operatorname{ord}_q(h^*F).\] G331's evaluation on
the DVR centre is \(f_F(v_q)=\operatorname{ord}_q(h^*F)/a_q\), including
the negative sign of a negative divisor. Pairing the independently
defined (5) gives the exact comparison
\[\int f_F\,d\mu_{\mathfrak a,R}=\deg_R(F).\tag{7}\] The normalized
degree is unchanged by any common primary-model pullback, since its
pullback to \(S\) is unchanged. Formula (6) also shows that the
independent Cartier zero-cycle degree on each such model is that same
degree. This proves actual model independence and exact Cartier
comparison for every real core presentation and every Cartier test in
this curve range.

For a primary ideal \(J\), apply one-variable multiplicity additivity
(complete Swanson--Huneke 11.2.3--4) to the finite normalization, whose
generic rank is one on every reduced component. Equivalently use n4's
already proved degree-one formula. The DVR Hilbert function gives
\(e_{S_q}(JS_q)=\operatorname{ord}_q(J)\), so
\[e_R(J)=\sum_q f_q\operatorname{ord}_q(J)=\int v(J)\,d\mu_{\mathfrak a,R}.\]
Thus (5) agrees with n3's zero-potential moment measure, but its Cartier
construction and proof (7) do not rely on assigning moments and then
renaming them intersections.

The bounded geometric function class on this finite set is every finite
nonpositive branch vector. Indeed each DVR's continuous cone consists of
nonpositive constants, obtained from powers of its parameter with real
coefficients; the semilocal class is the finite product. The
normalization fibres here are singletons, so G331's positivity descent
gives all such vectors downstairs. Its actual decreasing completion is
every extended nonpositive vector; bounded completed vectors are still
finite and hence already continuous. All-zero and mixed-zero profiles
are admissible as boundary functions. This says nothing about literal
logarithmic radial extensions at a node.

There are no potential variables in (5). Consequently two zero-tuples
give the identical measure on every open region of agreement; this is
actual open-set locality, albeit a vacuous zero-arity case. Every
bounded decreasing-net continuity statement for the zero-tuple is
constant equality of measures. It is not advertised as positive-arity
Bedford--Taylor locality. Finite continuous tests vary in the usual
finite-set norm, so their integrals are continuous as well. Unbounded or
bottom potential inputs are not being added: the tuple has zero entries.

For a primary \(\mathfrak b\) with
\(b_q=\operatorname{ord}_q(\mathfrak b)\), the exact change sends
\(v_q\) to \(\operatorname{ord}_q/b_q\) and multiplies its weight by
\(b_q/a_q\). Hence it is precisely (3), and normalization is (4) with
weights \(f_q\). On a smooth DVR over its own residue field, \(f_q=1\),
and (5) has mass \(a_q\) at \(\operatorname{ord}/a_q\); for the maximal
ideal the mass is one. This is the direct dimension-one normalized-local
reference convention. For a zero-dimensional reduced local germ, the
centred boundary is empty and the only measure is zero, without negative
arity.

A concrete residue-degree check prevents replacing \(f_q\) by one. Let
\(k\subset l\) be any finite field extension of degree \(d>1\), possibly
inseparable, and \[R=k+t\,l[[t]]\ \subset\ S=l[[t]].\] Put \(A=k[[t]]\).
Choosing a \(k\)-basis \(1,\alpha_2,\ldots,\alpha_d\) of \(l\) expresses
\(R=A\oplus\bigoplus_{i=2}^d tA\alpha_i\), a finite \(A\)-algebra, so it
is Noetherian and excellent. It is local with maximal ideal
\(\mathfrak m=tS\) and residue field \(k\): a nonzero constant in \(k\)
gives an inverse still having constant in \(k\). Its fraction field is
\(l((t))\) because \(t\alpha/t=\alpha\). Thus \(S\) is its finite
normalization with GENERIC degree one and RESIDUE degree \(d\). The
conductor is \(tS\): it contains that ideal, and a conductor unit would
force \(l=k\). For \(\mathfrak a=\mathfrak m\), (5) is
\(d\delta_{\operatorname{ord}_t}\). For \(J=t^rS\), \(J^N=t^{rN}S\) and
a \(k\)-basis of \(R/J^N\) consists of \(1\) and all \(\alpha_j t^i\),
\(1\le i<rN\), \(1\le j\le d\). Hence
\[\ell_R(R/J^N)=1+d(rN-1),\qquad e_R(J)=dr.\] This is exactly the moment
of (5) and the Cartier divisor degree. A generic-degree-only coefficient
would give \(r\), which is wrong. The argument does not assume
separability.

Artifacts: n5.md, bounded-evidence-input.json, exact\_diagnostics.py,
reports/exact-diagnostics.json. Rerun:
\texttt{python\ exact\_diagnostics.py};
\texttt{python\ build\_delivery.py}. The finite diagnostics check
declared arithmetic cases only. The universal curve and residue-degree
arguments above are mathematical proofs, not consequences of that
execution.

\section{n6. Precise degeneracies and a locality
diagnostic}\label{n6.-precise-degeneracies-and-a-locality-diagnostic}

Use G331 and n1--n5. These examples distinguish genuine obstructions
from missing proofs. They do not assert impossibility of the entire
reduced higher-dimensional P12.4 target.

First, nilpotent multiplicity cannot be reconstructed from the reduced
boundary and its gauge. Let \(R_1=k[[t]]\) and
\(R_2=k[[t]][\epsilon]/(\epsilon^2)\), both excellent local rings. The
ideal \((t)\) is primary in each, and their reduced gauges and centred
semivaluation boundaries are identical: nilpotents have infinite value.
Nevertheless
\[\ell_{R_1}(R_1/(t^N))=N,\qquad \ell_{R_2}(R_2/(t^N))=2N.\] The second
quotient has basis \(t^i,\epsilon t^i\) for \(0\le i<N\), while the
first has basis \(t^i\). Their Samuel multiplicities are one and two.
Thus no assignment depending ONLY on G331's reduced boundary and reduced
primary gauge can agree with both unreduced intersection multiplicities.
Extra generic-length data are necessary if that different target is
requested. Our measure construction deliberately uses \(R_{\rm red}\);
this is not an obstruction in the explicitly reduced scope.

Second, the degree-\(n\) convention has real lower-dimensional
degeneracy. Let
\[R=k[[x,y,z]]/(xy,xz)\cong k[[y,z]]\times_k k[[x]],\qquad n=2.\] It is
a reduced excellent complete local ring. The isomorphism sends a series
to its restrictions to \(x=0\) and \(y=z=0\), with equal constant terms;
the kernel is \((xy,xz)=(x)\cap(y,z)\). Its normalization is
\(S=k[[y,z]]\times k[[x]]\), finite because its quotient by \(R\) is one
copy of \(k\). It is normal on both component factors and has the same
total quotient fields. Its conductor is the pair of maximal ideals,
whose downstairs ideal is \(\mathfrak m=(x,y,z)\). Every centred
semivaluation has finite conductor value, so the normalized boundary is
the disjoint union of the FULL surface boundary and the curve point,
with no identified fibres. The pieces are clopen by G331.

Take \(\mathfrak a=\mathfrak m\) and \(I_N=(y,z,x^N)\). It is proper
primary. The rational profile \(u_N=-v(I_N)/N\) is constant \(-1/N\) on
the entire surface piece and \(-1\) at the curve point. These profiles
converge UNIFORMLY to the finite geometric \(C\) function \(u\) that is
zero on the surface and \(-1\) on the curve. This is uniform closure,
not an assumed increasing \(D\) closure. N4's top-degree normalization
formula gives \(e_R(I_N,\mathfrak m)=e_{k[[y,z]]}((y,z),(y,z))=1\); the
one-dimensional factor has no degree-two coefficient. The last equality
follows from the monomial basis count
\(\ell(k[[y,z]]/(y,z)^t)=t(t+1)/2\). Thus the moment measures of \(u_N\)
have total mass \(1/N\) and their weak limit is the zero measure. By n3,
\(\nu_{\mathfrak m}(u)=0\) even though \(u\ne0\). Formula (4) agrees:
the sole top-dimensional input pullback is zero. The lower curve point
was retained in the boundary and function class; it was not assigned a
vertical generic degree. To seek nonzero curve-local products on that
component would require the separate zero-potential
arity-one-dimensional convention, not silently adding it to a
one-potential dimension-two product.

Third, positivity, multilinearity and bounded continuity do not
logically imply actual locality. For a diagnostic independent of any
local ring, take the two-point compact discrete space \(K=\{x,y\}\), the
cone of all nonpositive finite functions, and the positive
measure-valued map \[F(u)=(-u(x)-u(y))\delta_x.\] It is additive,
nonnegative-real homogeneous, uniformly continuous, and continuous for
bounded decreasing nets. But \(u=(-1,-1)\) and \(w=(-1,-2)\) agree on
the open set \(U=\{x\}\), while \(F(u)|_U=2\delta_x\) and
\(F(w)|_U=3\delta_x\). This is an exact counterexample to the bare
implication from those properties to locality. It is NOT claimed to
satisfy the primary mixed-multiplicity kernel, to be a geometric
singular product, or to disprove locality of n3's measures. It explains
why a separate proof, tied to actual local intersections, is still
required.

The own program exact\_diagnostics.py enumerates only the finite bases
and rational gauge arithmetic stated in bounded-evidence-input.json,
with no supplied program imported or executed. The observed result was
PASS on 48 residue-degree basis cases (\(d=2,3\), \(r=1,2,3\),
\(N=1,\ldots,8\)), eight nilpotent basis cases (\(N=1,\ldots,8\)), and
27 two-branch gauge/test combinations, using exact fractions. It checks
the residue-degree length slopes, the nilpotent basis counts, and the
primary-gauge identity. This is exhaustive on those finite declared
arithmetic inputs only. Neither this evidence nor the toy map is a
substitute for higher-dimensional Cartier intersection comparison or
locality.

Rerun: \texttt{python\ exact\_diagnostics.py};
\texttt{python\ build\_delivery.py}. These diagnostics and their limits
are preserved as exact data, not a candidate solution claim.

\section{n7. Whole-original decision, unresolved proof coverage, and
deliverables}\label{n7.-whole-original-decision-unresolved-proof-coverage-and-deliverables}

Use qualified G331 and n1--n6. The ENTIRE original P12.4 is NOT solved
by this manuscript; candidate=false. This last note makes no
complete-solution claim. The positive results and gaps are separated
below so that a reviewer can use the partial mathematics without
importing a false product theorem.

What is established in the draft proof chain is: an unconditional
mixed-multiplicity kernel at every excellent reduced local ring, a
unique finite positive moment-measure extension to every real \(C\)
tuple and every bounded input in the ACTUAL decreasing completion \(D\),
a weak descending-net theorem whose specified limit belongs to that
domain, exact primary-gauge and normalization identities for that
construction, and identification with the exact published BFJ operator
in its smooth affine CLOSED-germ algebraically closed
characteristic-zero range. Every semivaluation support and full
conductor fibre is retained. Inputs with zero or mixed-zero coordinates
are not excluded; unbounded and bottom inputs are outside this bounded
construction. Non-equidimensional rings use the explicitly stated
degree-\(n\) convention and \(n-1\) potential slots.

There is also an actual geometric result for all excellent reduced
curves: the zero-potential reference measure (5), independent Cartier
degree comparison (6)--(7) on actual primary models for real
presentations, model independence, primary-gauge and normalization rules
with actual residue-degree weights, and open-set locality and bounded
continuity for that zero-arity operator. Dimension zero has the
empty-boundary zero measure. N6 supplies a genuine nilpotent-data
obstruction and an exact non-equidimensional degeneracy, with their
limited conclusions stated.

Two essential higher-dimensional geometric obligations remain
unresolved:

\begin{enumerate}
\def\labelenumi{\arabic{enumi}.}
\item
  For every actual negative curve-nef REAL Cartier presentation on an
  actual primary blowup, prove that the scalar kernel extension is its
  independent local Cartier intersection, with all nonnormal,
  non-Cohen--Macaulay, finite-residue-field and non-equidimensional
  issues handled. In particular, the real \(C\) comparison supplied by
  G331 is a FUNCTION approximation theorem, not an
  INTERSECTION-continuity theorem. Knowing that section ideals
  approximate the evaluation uniformly does not prove their
  multiplicities converge to an independently defined
  Euler-characteristic intersection. No arbitrary-model cofinality is
  available. A correct projection and independent-intersection argument
  on these primary models must be supplied.
\item
  Prove actual locality on open regions: if each pair \(u_i,w_i\) of
  bounded geometric inputs agrees on an open
  \(U\subset X_{\mathfrak a}\), prove equality of the restrictions of
  their mixed measures to \(U\). Global length polynomials, positivity
  on tests, uniqueness of moment representation, and weak limits do not
  establish this. The finite toy diagnostic in n6 shows the general
  logical implication is false. Even a localized atom formula for core
  inputs would require a further passage to arbitrary bounded completed
  inputs; no such passage has been proved here. BFJ Section 4 supplies
  its exact smooth construction/comparison/continuity, not a singular
  arbitrary-characteristic locality grant.
\end{enumerate}

Therefore n3's measures are called moment measures, not claimed to be
P12.4's actual geometric products. Their algebraic presentation
independence must not be relabelled as geometric model-independent
intersection theory. The exact full original is left open in dimension
at least two. No general obstruction to the reduced target has been
proved; missing proof is not impossibility. The nilpotent example
obstructs only a boundary-only unreduced multiplicity reconstruction.
P12.5 and general P12.6 remain separate.

The fresh deliverables are P12.4-paper.tex/PDF, argument.md,
source-application-ledger.md, exact-scope.md, whole-proof-coverage.md,
original-scope-and-noncircularity-decision.md, native-notes.json, n1.md
through n7.md, bounded-evidence-input.json, exact\_diagnostics.py,
audit\_inputs.py, build\_delivery.py, and final\_audit.py. Exact bound
mathematical/primary source bodies and manifests are retained under
inputs; reports retain their URLs, original/local paths, hashes,
declared reading units, environment details, finite evidence and build
results. The incoming qualified G331 source, mandatory annotation and
source ledger are copied unchanged. No attached/source/ROOT program was
executed or imported. No active private campaign output or hot feedback
was read.

Rerun from this directory: \texttt{python\ audit\_inputs.py};
\texttt{python\ exact\_diagnostics.py};
\texttt{python\ build\_delivery.py}; \texttt{python\ final\_audit.py}.
Compilation uses pdflatex with no shell escape. All scripts are freshly
owned and execute only the declared small local tasks. No supervised
heavy-compute service was supplied or needed. No agents,
authentication/runtime/environment/HOME/CODEX\_HOME/proxy/CA changes,
publication, catalogue changes, or protected editor interactions were
used. No tool output is treated as mathematical verification. The
caller's requested model name is not an independently measured execution
attestation.

The final audit checks file/support/hash consistency only. Every paper
claim remains an unverified mathematical draft. The candidate flag
remains false specifically because independent higher-dimensional
Cartier comparison and locality are not proved; artifact completeness
and a successful build cannot change that decision.

\end{document}
